\begingroup
\color{black}
\hypersetup{citecolor=black,linkcolor=black,urlcolor=black}
\section{Conditions for Allocation Incentives}
\label{app:operational-incentives}

The operational term of \RSCI decreases when a workload uses less energy or runs during cleaner hours or in cleaner regions.
By construction, total attributed emissions equal the provider's reported total.
The analysis below holds that report fixed to examine how allocation affects an individual workload's score, giving the incentive conditions discussed alongside Table~\ref{tab:properties}.

\parag{Setup}
Consider one reporting slice and suppress its index $s$.
Let $x\geq0$ be the workload's modeled operational emissions and $X>0$ the total across all workloads in the slice.
Hold reported emissions $C^{\downarrow k}\geq0$, scope fractions $\beta^k\in[0,1]$, functional output $R_i>0$, other workloads' operational emissions, and capacity shares fixed.
Write
\begin{equation}
A=\sum_k\beta^k C^{\downarrow k},\qquad
B=\sum_k(1-\beta^k)C^{\downarrow k},\qquad
b=\beta^{\mathrm{S2}},\qquad v=w^{\mathrm{cap}}_{i,s}\in[0,1].
\end{equation}
The energy and capacity residuals are $A-bX$ and $B-(1-b)X$, respectively.
With energy weights proportional to operational emissions, this workload's attributed emissions are
\begin{equation}
\begin{aligned}
F(x)&=x+\frac{x}{X}(A-bX)+v\bigl(B-(1-b)X\bigr)\\
    &=(1-b)x+A\frac{x}{X}+vB-v(1-b)X.
\end{aligned}
\end{equation}
Dividing by the fixed $R_i$ gives this slice's contribution to the workload's score.

\parag{Operational improvements}
Because other workloads' emissions are fixed, $dX/dx=1$, giving
\begin{equation}
\frac{dF}{dx}=(1-b)(1-v)+A\frac{X-x}{X^2}\geq0.
\label{eq:operational-incentive-derivative}
\end{equation}
\emph{Under these assumptions, reducing operational emissions cannot increase the reconciled score} while $X$ remains positive.
Both the residuals and the energy weights are recomputed after the change.
The derivative is strictly positive if either $b<1$ and $v<1$, or $A>0$ and other workloads have positive operational emissions.
In particular:
\begin{itemize}
\setlength{\itemsep}{1pt}
\item If $\beta^{\mathrm{S2}}<1$, a workload with less than the full capacity share receives a strictly lower score when its operational emissions decrease.
\item If $\beta^{\mathrm{S2}}=1$, positive reported Scope~2 emissions and positive operational emissions from other workloads suffice for a strict decrease. Energy allocations from other scopes can also make $A>0$.
\end{itemize}
These conditions cover energy savings and moving work to cleaner hours when operational emissions and energy weights use the time-varying carbon intensity, and capacity shares remain fixed.
The result extends across slices if the stated assumptions hold in each slice and operational emissions do not increase in any affected slice.
When work moves between slices, lower total operational emissions alone do not guarantee a lower score. The combined attributions must be compared.

\parag{Capacity allocation}
Now hold all operational emissions, reports, scope fractions, and output fixed, and vary the workload's capacity share $v$.
Its effect on the score is
\begin{equation}
\frac{\partial (F/R_i)}{\partial v}=\frac{B-(1-b)X}{R_i}=\frac{\Delta_s^{\mathrm{cap}}}{R_i}.
\end{equation}
\emph{A smaller capacity share lowers the score for a positive capacity residual, but raises it for a negative residual.}
If the capacity residual is zero, changing the capacity share has no effect on the score.
Setting every $\beta^k=1$ removes the capacity component.

\parag{Limits}
For a sole workload, $x=X$ and $v=1$, so $F=A+B=\sum_k C^{\downarrow k}$.
\emph{Its total attribution follows the provider's reported total.}
It therefore remains unchanged in the fixed-report analysis above.
These results concern attribution and do not establish physical emissions savings.
Actions that also change capacity shares, provider reports, scope fractions, other workloads' emissions, or functional output must be assessed with those changes included.
\endgroup
